% ===========================================================================
% Gesamtskript VWL -- Musterlösungen Lektion 1: Grundbegriffe & Preisbildung
% (Fragment, wird per \input eingebunden; Aufgaben in aufgaben-l1.tex)
% ===========================================================================

% ---------------------------------------------------------------------------
\begin{loesung}[title={Lösung Aufgabe~1.1 --- Wirtschaften und ökonomisches Prinzip}]
\textbf{a) Bedürfnis, Bedarf, Nachfrage} \hfill (3~P.: je Begriff 1~P.)
\begin{itemize}
  \item \emph{Bedürfnis:} Frau Jansen empfindet einen Mangel --- sie möchte ihre Kinder
    ohne Auto zur Kita bringen und wünscht sich dafür ein geeignetes Verkehrsmittel.
  \item \emph{Bedarf:} Da sie 4.500~€ angespart hat, ist ihr Bedürfnis mit
    \emph{Kaufkraft} ausgestattet; aus dem Bedürfnis wird Bedarf.
  \item \emph{Nachfrage:} Mit der Bestellung im Onlineshop wird der Bedarf am Markt
    \emph{wirksam} --- es entsteht Nachfrage nach einem Lastenrad.
\end{itemize}

\textbf{b) Zuordnung zum ökonomischen Prinzip} \hfill (6~P.: je Aussage 1~P. Zuordnung +
1~P. Begründung)
\begin{itemize}
  \item \emph{Aussage (1) --- Minimalprinzip:} Das Ziel ist vorgegeben (2.400 Rahmen); es
    soll mit möglichst geringem Mitteleinsatz (niedrigste Beschaffungskosten) erreicht
    werden.
  \item \emph{Aussage (2) --- Maximalprinzip:} Die Mittel sind vorgegeben (Messebudget
    18.000~€); mit ihnen soll ein möglichst großer Erfolg (Zahl der Händlerkontakte)
    erzielt werden.
  \item \emph{Aussage (3) --- kein ökonomisches Prinzip:} Mittel und Ziel werden
    \emph{gleichzeitig} optimiert („geringstmögliche Kosten`` \emph{und}
    „größtmögliche Zahl``). Das ist logisch nicht lösbar: Die geringstmöglichen Kosten
    wären null, dann würde nichts produziert. Eine der beiden Größen muss fest vorgegeben
    sein.
\end{itemize}

\textbf{c) Produktionsfaktoren} \hfill (2~P.: je Zuordnung 0,5~P.)

Montageteam $\rightarrow$ \emph{Arbeit}; Gewerbegrundstück $\rightarrow$ \emph{Boden}
(Standort); Lackierroboter $\rightarrow$ \emph{Kapital} (produziertes Produktionsmittel);
patentiertes Klebeverfahren $\rightarrow$ \emph{Wissen (Know-how)}.
\end{loesung}

% ---------------------------------------------------------------------------
\begin{loesung}[title={Lösung Aufgabe~1.2 --- Marktarten und Marktformen}]
\textbf{a) Marktarten} \hfill (4~P.: je Vorgang 1~P.)
\begin{enumerate}[label=(\arabic*)]
  \item Stellenausschreibung $\rightarrow$ \emph{Faktormarkt} (Arbeitsmarkt).
  \item Anmietung der Lagerhalle $\rightarrow$ \emph{Faktormarkt} (Immobilienmarkt,
    Faktor Boden).
  \item Investitionskredit $\rightarrow$ \emph{Geld- und Kapitalmarkt} (Finanzmarkt,
    Kreditmarkt).
  \item Verkauf von Kartoffeln $\rightarrow$ \emph{Gütermarkt} (Konsumgütermarkt).
\end{enumerate}

\textbf{b) Marktformen} \hfill (6~P.: je Situation 1~P. Marktform + 1~P. Begründung)
\begin{enumerate}[label=(\arabic*)]
  \item \emph{Nachfrageoligopol:} Viele Anbieter (rund 1.800 Landwirte) stehen
    \emph{wenigen} Nachfragern (zwei Zuckerfabriken) gegenüber. Die Marktmacht liegt bei
    den Nachfragern.
  \item \emph{Angebotsmonopol:} Ein einziger Anbieter (die Reederei) steht vielen
    Nachfragern (Fahrgästen) gegenüber; Ausweichmöglichkeiten fehlen weitgehend.
  \item \emph{Angebotsoligopol:} Wenige Anbieter (internationale Chemiekonzerne) mit
    großen Marktanteilen stehen vielen Nachfragern (Händler, Landwirte) gegenüber.
\end{enumerate}

\textbf{c) Beurteilung der Aussage} \hfill (2~P.)

Die Aussage ist \emph{falsch}. Viele Anbieter und viele Nachfrager (Polypol) sind nur
\emph{eine} von fünf Bedingungen des vollkommenen Marktes. Auf dem Wochenmarkt fehlen
weitere Bedingungen, z.\,B.: Die Güter sind nicht homogen (Bio- und konventionelle Ware,
unterschiedliche Sorten und Qualität), Kunden haben persönliche Präferenzen
(Stammhändler) und räumliche Präferenzen (Stand am Eingang), und es besteht keine
vollständige Markttransparenz. Es handelt sich daher um einen \emph{unvollkommenen} Markt
(unvollkommenes Polypol).
\end{loesung}

% ---------------------------------------------------------------------------
\begin{loesung}[title={Lösung Aufgabe~1.3 --- Preiselastizität und Umsatzwirkung}]
\textbf{a) Direkte Preiselastizität} \hfill (4~P.: je Produkt 2~P.)

Formel: $\displaystyle \varepsilon = \frac{\text{Mengenänderung in \%}}{\text{Preisänderung in \%}}$
\quad (Änderungen jeweils bezogen auf den Ausgangswert)

\smallskip
\emph{Espresso „Crema Oro``:}
\begin{itemize}
  \item Preisänderung: $\dfrac{15{,}40~\text{€} - 14{,}00~\text{€}}{14{,}00~\text{€}}
    = \dfrac{1{,}40~\text{€}}{14{,}00~\text{€}} = +10\,\%$
  \item Mengenänderung: $\dfrac{6.800 - 8.000}{8.000} = \dfrac{-1.200~\text{Pck.}}{8.000~\text{Pck.}} = -15\,\%$
  \item $\varepsilon = \dfrac{-15\,\%}{+10\,\%} = \mathbf{-1{,}5}$ \quad $\Rightarrow\ |\varepsilon| = 1{,}5$
\end{itemize}

\emph{Filterkaffee „Hausmischung``:}
\begin{itemize}
  \item Preisänderung: $\dfrac{6{,}60~\text{€} - 6{,}00~\text{€}}{6{,}00~\text{€}}
    = \dfrac{0{,}60~\text{€}}{6{,}00~\text{€}} = +10\,\%$
  \item Mengenänderung: $\dfrac{19.200 - 20.000}{20.000} = \dfrac{-800~\text{Pck.}}{20.000~\text{Pck.}} = -4\,\%$
  \item $\varepsilon = \dfrac{-4\,\%}{+10\,\%} = \mathbf{-0{,}4}$ \quad $\Rightarrow\ |\varepsilon| = 0{,}4$
\end{itemize}

\textbf{b) Interpretation} \hfill (4~P.: je Produkt 1~P. Einordnung + 1~P. Begründung)
\begin{itemize}
  \item \emph{Espresso:} $|\varepsilon| = 1{,}5 > 1$ $\Rightarrow$ \emph{elastische}
    Nachfrage. Die Menge sinkt prozentual stärker (15~\%) als der Preis steigt (10~\%).
    Begründung: Premium-Espresso ist kein Grundbedarf; es gibt viele Ausweichmöglichkeiten
    (andere Marken, Discounter-Espresso, Kapselsysteme) bzw. die Kunden können auf das
    Produkt verzichten.
  \item \emph{Filterkaffee:} $|\varepsilon| = 0{,}4 < 1$ $\Rightarrow$ \emph{unelastische}
    Nachfrage. Die Menge sinkt prozentual schwächer (4~\%) als der Preis steigt (10~\%).
    Begründung: Filterkaffee ist für viele Haushalte ein Alltagsgut mit festen
    Konsumgewohnheiten; die Kunden sind an „ihre`` Hausmarke gebunden und reagieren auf
    Preisänderungen kaum.
\end{itemize}
Das negative Vorzeichen ist bei der direkten Preiselastizität normal (Preis steigt,
Menge sinkt); maßgeblich für die Einordnung ist der Betrag.

\textbf{c) Umsatzwirkung} \hfill (4~P.: je Produkt 2~P.)

\begin{center}
\small
\begin{tabular}{l r r r r}
\toprule
\rowcolor{bfwTabHead}
\textbf{Produkt} & \textbf{Umsatz vorher} & \textbf{Umsatz nachher} & \textbf{Änderung in €} & \textbf{Änderung in \%}\\
\midrule
Espresso & $14{,}00 \cdot 8.000 = 112.000$~€ & $15{,}40 \cdot 6.800 = 104.720$~€ & $-7.280$~€ & $-6{,}5\,\%$\\
Filterkaffee & $6{,}00 \cdot 20.000 = 120.000$~€ & $6{,}60 \cdot 19.200 = 126.720$~€ & $+6.720$~€ & $+5{,}6\,\%$\\
\bottomrule
\end{tabular}
\end{center}
Prozentwerte: $-7.280 \div 112.000 = -0{,}065 = -6{,}5\,\%$;
$\ +6.720 \div 120.000 = +0{,}056 = +5{,}6\,\%$.\\
Probe: $1{,}10 \cdot 0{,}85 = 0{,}935$ ($-6{,}5\,\%$) bzw. $1{,}10 \cdot 0{,}96 = 1{,}056$
($+5{,}6\,\%$). Ergebnis: Beim elastischen Gut sinkt, beim unelastischen Gut steigt der
Umsatz.

\textbf{d) Beurteilung des Vorschlags} \hfill (4~P.)
\begin{itemize}
  \item \emph{Espresso:} Die Preiserhöhung ist aus Umsatzsicht \emph{nicht} zu empfehlen,
    da die elastische Nachfrage zu einem Umsatzrückgang von 7.280~€ je Quartal führt.
    Zudem droht der dauerhafte Verlust von Kunden an die Konkurrenz.
  \item \emph{Filterkaffee:} Die Preiserhöhung ist zu empfehlen; der Umsatz steigt um
    6.720~€ je Quartal.
  \item \emph{Gesamtbetrachtung:} Würden beide Erhöhungen umgesetzt, sänke der
    Gesamtumsatz um $7.280 - 6.720 = 560$~€ je Quartal. Sinnvoll ist daher eine
    \emph{differenzierte} Preispolitik (nur Filterkaffee erhöhen).
  \item \emph{Einschränkungen:} Umsatz ist nicht gleich Gewinn --- durch die geringere
    Absatzmenge sinken auch die variablen Kosten; die Entscheidung sollte daher auch über
    den Deckungsbeitrag geprüft werden. Außerdem gelten die Testmarktdaten nur
    \emph{ceteris paribus} und kurzfristig; bundesweit und langfristig (Reaktion der
    Wettbewerber, Gewöhnung der Kunden) können die Elastizitäten abweichen.
\end{itemize}
\end{loesung}

% ---------------------------------------------------------------------------
\begin{loesung}[title={Lösung Aufgabe~1.4 --- Marktgleichgewicht aus Angebots- und Nachfragefunktion}]
\textbf{a) Gleichgewicht} \hfill (4~P.: $p^*$ 2~P., $x^*$ 1~P., Umsatz 1~P.)

Im Gleichgewicht gilt: angebotene Menge = nachgefragte Menge, also $x_A = x_N$.
\[
60\,p - 300 = 1.200 - 40\,p
\quad\Rightarrow\quad 100\,p = 1.500
\quad\Rightarrow\quad p^* = \mathbf{15~\text{€ je Sack}}
\]
Einsetzen: $x_N = 1.200 - 40 \cdot 15 = 1.200 - 600 = \mathbf{600~\text{Säcke pro Tag}}$\\
Probe: $x_A = 60 \cdot 15 - 300 = 900 - 300 = 600$~Säcke (stimmt überein)\\
Umsatz im Gleichgewicht: $15~\text{€/Sack} \cdot 600~\text{Säcke} = \mathbf{9.000~\text{€ pro Tag}}$.

\textbf{b) Überhänge bei 12~€ und 20~€} \hfill (6~P.: je Preis 1~P. Mengen, 1~P. Überhang,
1~P. Marktsituation/Reaktion)

\begin{center}
\small
\begin{tabular}{c c c l}
\toprule
\rowcolor{bfwTabHead}
\textbf{Preis} & \textbf{Nachfrage $x_N$} & \textbf{Angebot $x_A$} & \textbf{Überhang}\\
\midrule
12~€ & $1.200 - 40 \cdot 12 = 720$ & $60 \cdot 12 - 300 = 420$ & Nachfrageüberhang $720 - 420 = 300$~Säcke\\
20~€ & $1.200 - 40 \cdot 20 = 400$ & $60 \cdot 20 - 300 = 900$ & Angebotsüberhang $900 - 400 = 500$~Säcke\\
\bottomrule
\end{tabular}
\end{center}
\begin{itemize}
  \item \emph{Preis 12~€ (unter $p^*$):} Nachfrageüberhang (Angebotsdefizit) von
    300~Säcken pro Tag $\Rightarrow$ \emph{Verkäufermarkt}. Nicht alle Kunden werden
    beliefert; die Anbieter können höhere Preise durchsetzen. Mit steigendem Preis
    weiten die Anbieter ihre Menge aus, die nachgefragte Menge sinkt --- der Preis
    steigt in Richtung 15~€.
  \item \emph{Preis 20~€ (über $p^*$):} Angebotsüberhang (Nachfragedefizit) von
    500~Säcken pro Tag $\Rightarrow$ \emph{Käufermarkt}. Die Anbieter bleiben auf Ware
    sitzen und senken die Preise, um abzusetzen. Mit sinkendem Preis steigt die
    nachgefragte Menge, die angebotene Menge sinkt --- der Preis fällt in Richtung 15~€.
\end{itemize}
Ergänzend: Abgesetzt wird jeweils nur die Menge der „kürzeren`` Marktseite (bei 12~€:
420~Säcke, bei 20~€: 400~Säcke). Nur im Gleichgewicht wird mit 600~Säcken die größtmögliche
Menge umgesetzt.

\textbf{c) Preis-Mengen-Diagramm} \hfill (4~P.: Achsen 1~P., beide Kurven 1~P.,
Gleichgewicht 1~P., Überhang 1~P.)

Hilfspunkte: Nachfrage $p = 30 - x/40$ (Prohibitivpreis 30~€ bei $x = 0$; Sättigungsmenge
1.200~Säcke bei $p = 0$). Angebot $p = 5 + x/60$ (Mindestpreis 5~€ bei $x = 0$; bei
$x = 1.200$ ist $p = 25$~€).

\begin{center}
\begin{tikzpicture}[x=1cm, y=1cm, >=Stealth, font=\small]
  % Renten (für Teilaufgabe d)
  \fill[bfwAccent!25] (0,6) -- (0,3) -- (5,3) -- cycle;
  \fill[bfwPrimary!20] (0,1) -- (0,3) -- (5,3) -- cycle;
  % Achsen
  \draw[->, thick, bfwInk] (0,0) -- (10.8,0) node[right] {$x$};
  \draw[->, thick, bfwInk] (0,0) -- (0,6.8) node[above] {Preis $p$ (€/Sack)};
  \node[below, font=\footnotesize] at (5,-0.5) {Menge $x$ (Säcke pro Tag)};
  \foreach \x in {200,400,600,800,1000,1200}
    \draw (\x/120,0.08) -- (\x/120,-0.08) node[below] {\footnotesize\x};
  \foreach \p in {5,10,15,20,25,30}
    \draw (0.08,\p/5) -- (-0.08,\p/5) node[left] {\footnotesize\p};
  % Kurven
  \draw[very thick, bfwPrimaryDark] (0,6) -- (10,0) node[pos=0.97, above right] {$N$};
  \draw[very thick, bfwWarn] (0,1) -- (10,5) node[right] {$A$};
  % Gleichgewicht
  \draw[dashed, bfwInkSoft] (0,3) -- (5,3) -- (5,0);
  \fill[bfwInk] (5,3) circle (2.5pt);
  \node[below right, bfwInk] at (5.1,2.95) {GG ($x^*=600$; $p^*=15$)};
  % Überhang bei 20 Euro
  \draw[dashed, bfwInkSoft] (0,4) -- (7.5,4);
  \fill[bfwPrimaryDark] (10/3,4) circle (2pt);
  \fill[bfwWarn] (7.5,4) circle (2pt);
  \draw[<->, thick, bfwAccent!85!black] (10/3,4.35) -- (7.5,4.35)
    node[midway, above] {Angebotsüberhang 500};
  % Beschriftung Renten
  \node[font=\footnotesize] at (1.4,3.7) {KR};
  \node[font=\footnotesize] at (1.4,2.35) {PR};
\end{tikzpicture}
\end{center}

\textbf{d) Konsumentenrente} \hfill (3~P.: Rechnung 2~P., Erläuterung 1~P.)

Die Konsumentenrente ist das Dreieck zwischen Nachfragekurve und Gleichgewichtspreis
(Fläche KR im Diagramm):
\[
KR = \tfrac{1}{2} \cdot x^* \cdot (\text{Prohibitivpreis} - p^*)
   = \tfrac{1}{2} \cdot 600~\text{Säcke} \cdot (30~\text{€} - 15~\text{€})
   = \mathbf{4.500~\text{€ pro Tag}}
\]
\emph{Erläuterung:} Viele Käufer wären bereit gewesen, mehr als 15~€ je Sack zu zahlen.
Da sie nur den einheitlichen Marktpreis von 15~€ zahlen, entsteht ihnen ein
Nutzengewinn (eine „Ersparnis``) von insgesamt 4.500~€ pro Tag.
(Zum Vergleich: Produzentenrente $= \tfrac{1}{2} \cdot 600 \cdot (15 - 5) = 3.000$~€ pro Tag.)
\end{loesung}

% ---------------------------------------------------------------------------
\begin{loesung}[title={Lösung Aufgabe~1.5 --- Kurvenverschiebungen und Kreuzpreiselastizität}]
\textbf{a) Bewegung oder Verschiebung} \hfill (8~P.: je Entwicklung 2~P.)
\begin{enumerate}[label=(\arabic*)]
  \item \emph{Hitzesommer:} Die Bedürfnisstruktur bzw. Zukunftserwartung ändert sich,
    nicht der Preis $\Rightarrow$ \emph{Verschiebung der Nachfragekurve nach rechts}
    (Nachfrageerhöhung): Bei jedem Preis wird mehr Limonade nachgefragt.
  \item \emph{Zuckerpreis $+40\,\%$:} Die Kosten eines Produktionsfaktors (Rohstoff)
    steigen $\Rightarrow$ \emph{Verschiebung der Angebotskurve nach links}
    (Angebotssenkung): Bei jedem Preis wird weniger angeboten.
  \item \emph{Preis der Limonade sinkt:} Es ändert sich der Preis des Gutes selbst
    $\Rightarrow$ \emph{keine Verschiebung}, sondern \emph{Bewegung auf den Kurven}:
    auf der Nachfragekurve nach rechts unten (nachgefragte Menge steigt), auf der
    Angebotskurve nach links unten (angebotene Menge sinkt).
  \item \emph{Neue Abfülltechnik:} Technischer Fortschritt erhöht die Produktivität und
    senkt die Stückkosten $\Rightarrow$ \emph{Verschiebung der Angebotskurve nach rechts}
    (Angebotserhöhung).
\end{enumerate}

\textbf{b) Gleichzeitiger Eintritt von (1) und (2)} \hfill (4~P.)

Die Nachfragekurve verschiebt sich nach rechts, die Angebotskurve nach links. Beide
Verschiebungen wirken \emph{preissteigernd} $\Rightarrow$ Der Gleichgewichtspreis
\emph{steigt eindeutig}. Für die Gleichgewichtsmenge wirken die Effekte
\emph{gegenläufig}: Die Nachfrageerhöhung erhöht, die Angebotssenkung verringert die
Menge. Die Wirkung auf die Gleichgewichtsmenge ist daher \emph{unbestimmt}; sie hängt
davon ab, welche Verschiebung stärker ausfällt (überwiegt die Nachfrageerhöhung, steigt
die Menge; überwiegt die Angebotssenkung, sinkt sie).

\textbf{c) Kreuzpreiselastizität} \hfill (4~P.: Mengenänderung 1~P., Elastizität 2~P.,
Güterbeziehung 1~P.)

Mengenänderung Limonade: $\dfrac{11.750 - 12.500}{12.500} = \dfrac{-750~\text{Kisten}}{12.500~\text{Kisten}} = -6\,\%$
\[
\varepsilon_K = \frac{\text{Mengenänderung Limonade in \%}}{\text{Preisänderung Eistee in \%}}
= \frac{-6\,\%}{-20\,\%} = \mathbf{+0{,}3}
\]
Das \emph{positive} Vorzeichen zeigt: Preis des Eistees und Menge der Limonade bewegen
sich gleichgerichtet (Eistee billiger $\rightarrow$ weniger Limonade). Eistee und Limonade
sind \emph{Substitutionsgüter}. Da $|\varepsilon_K| = 0{,}3$ relativ klein ist, ist die
Substitutionsbeziehung nur schwach ausgeprägt.
\end{loesung}

% ---------------------------------------------------------------------------
\begin{loesung}[title={Lösung Aufgabe~1.6 --- Funktionen des Marktpreises}]
\textbf{a) Funktionen des Marktpreises} \hfill (6~P.: je Funktion 1~P. Erläuterung +
1~P. Situationsbezug; drei der folgenden)
\begin{itemize}
  \item \emph{Signalfunktion (Informationsfunktion):} Der Preis zeigt die Knappheit eines
    Gutes an. Der Anstieg von 420~€ auf 680~€ je m³ signalisiert allen Marktteilnehmern,
    dass Konstruktionsvollholz knapp geworden ist.
  \item \emph{Allokationsfunktion (Lenkungsfunktion):} Hohe Preise lenken
    Produktionsfaktoren in die Herstellung des knappen Gutes. Die Sägewerke setzen mehr
    Arbeit ein (Sonderschichten), ein Investor lenkt Kapital in ein neues Sägewerk.
  \item \emph{Selektionsfunktion (Sanktionsfunktion):} Marktteilnehmer, die den Preis
    nicht zahlen bzw. nicht weitergeben können, scheiden aus. Kleine Zimmereien, die die
    höheren Einkaufspreise nicht an ihre Kunden weitergeben können, geben auf.
  \item \emph{Ausgleichsfunktion:} Der höhere Preis bringt Angebots- und Nachfragemenge
    wieder zum Ausgleich: Die Nachfrager schränken ihre Käufe ein (z.\,B. Ausweichen auf
    andere Baustoffe), die Anbieter weiten ihre Menge aus, bis der Markt geräumt ist.
\end{itemize}

\textbf{b) Marktsituation und Preismechanismus} \hfill (4~P.)

Durch den Nachfrageanstieg hat sich die Nachfragekurve nach rechts verschoben. Beim alten
Preis von 420~€ je m³ übersteigt die nachgefragte Menge die angebotene Menge: Es besteht
ein \emph{Nachfrageüberhang} (Angebotsdefizit), also ein \emph{Verkäufermarkt} (1~P.).
Nicht alle Kaufwünsche können erfüllt werden; die Nachfrager überbieten sich, die
Anbieter setzen höhere Preise durch (1~P.). Mit steigendem Preis weiten die Sägewerke
ihre angebotene Menge aus (Bewegung auf der Angebotskurve) und ein Teil der Nachfrager
verzichtet oder weicht aus (Bewegung auf der neuen Nachfragekurve) (1~P.). Der Prozess
endet im neuen Gleichgewicht mit \emph{höherem} Preis und \emph{höherer} Menge als vor dem
Nachfrageanstieg (1~P.).
\end{loesung}

% ---------------------------------------------------------------------------
\begin{loesung}[title={Lösung Aufgabe~1.7 --- Monopolistischer Absatzbereich im unvollkommenen Polypol}]
\textbf{a) Marktform und Preisspielraum} \hfill (4~P.)

Es liegt ein \emph{unvollkommenes Polypol} vor: viele Anbieter, viele Nachfrager
(Polypol, 1~P.), aber ein unvollkommener Markt, weil die Güter \emph{heterogen} sind
(unterschiedliche Marken, Rezepturen, Verpackungen) und die Kunden \emph{Präferenzen}
haben (1~P.). Durch \emph{Produktdifferenzierung} --- hier das Steingut-Glas, das
traditionelle Etikett, die Marke „Altmühltaler`` und die besondere Rezeptur --- bindet
die Brinkhaus GmbH Stammkunden an sich (1~P.). Daraus entsteht ein
\emph{monopolistischer Absatzbereich} (zwischen 2,20~€ und 3,10~€): Innerhalb dieses
Bereichs kann das Unternehmen seinen Preis wie ein Monopolist weitgehend frei festlegen,
ohne größere Nachfrageeinbußen befürchten zu müssen (1~P.).

\textbf{b) Beurteilung der Preiserhöhungen} \hfill (4~P.: je Fall 2~P.)
\begin{enumerate}[label=(\arabic*)]
  \item \emph{2,95~€:} Der Preis liegt noch \emph{innerhalb} des monopolistischen
    Absatzbereichs (unter 3,10~€). Die Stammkunden bleiben weitgehend treu; der Absatz
    geht nur geringfügig zurück. Die Nachfrage reagiert hier unelastisch, der Umsatz
    dürfte steigen.
  \item \emph{3,40~€:} Der Preis liegt \emph{oberhalb} der Obergrenze von 3,10~€. Die
    Präferenzen reichen nicht mehr aus; viele Kunden wandern zu Konkurrenzmarken ab. Die
    Nachfrage reagiert hier stark elastisch, Absatz und Umsatz brechen deutlich ein.
    Die Erhöhung ist nicht zu empfehlen.
\end{enumerate}

\textbf{c) Maßnahmen zur Vergrößerung des Spielraums} \hfill (2~P.: je Maßnahme 1~P.)

Zum Beispiel: Ausbau der Markenbekanntheit durch Werbung; Qualitätsverbesserung bzw.
Betonung der Herkunft (Regionalität, Bio-Siegel); weitere Differenzierung durch
Sondereditionen oder Geschenkverpackungen; Kundenbindungsmaßnahmen und Service;
Vertrieb über exklusive Kanäle (Feinkosthandel).
\end{loesung}

% ---------------------------------------------------------------------------
\begin{loesung}[title={Lösung Aufgabe~1.8 --- Preisbildung im Oligopol und im Monopol}]
\textbf{a) Marktform} \hfill (2~P.)

\emph{Angebotsoligopol} (1~P.): Wenige Anbieter (drei Zementwerke) mit großen
Marktanteilen stehen vielen Nachfragern (mehrere Hundert Abnehmer) gegenüber (1~P.).

\textbf{b) Typische Verhaltensweisen} \hfill (4~P.: je Verhaltensweise 2~P.; zwei der
folgenden)
\begin{itemize}
  \item \emph{Parallelverhalten / Preisruhe:} Die Anbieter wissen, dass jede eigene
    Preisänderung die Konkurrenten unmittelbar trifft und eine Reaktion auslöst. Sie
    verhalten sich deshalb ohne ausdrückliche Absprache gleichgerichtet und lassen die
    Preise über längere Zeit unverändert (Preisstarre); Änderungen erfolgen im
    Gleichschritt.
  \item \emph{Preisführerschaft:} Ein (meist großer) Anbieter erhöht als Erster die
    Preise, z.\,B. wegen gestiegener Energiekosten; die übrigen Werke ziehen nach.
  \item \emph{Verlagerung des Wettbewerbs:} Statt über den Preis konkurrieren die Werke
    über Qualität, Lieferzuverlässigkeit, Service und Konditionen.
  \item \emph{Marktverdrängungspolitik:} Ein Anbieter versucht, mit aggressiven
    Preissenkungen Marktanteile zu gewinnen bzw. Konkurrenten zu verdrängen (selten).
\end{itemize}

\textbf{c) Beurteilung der Preissenkung um 10\,\%} \hfill (3~P.)

Die Erfolgsaussichten sind gering. Im Oligopol besteht eine enge
\emph{Reaktionsverbundenheit}: Eine Preissenkung von 10~\% zieht spürbar Kunden der
beiden Konkurrenten ab, die deshalb mit hoher Wahrscheinlichkeit sofort nachziehen
(1~P.). Die Marktanteile bleiben dann annähernd gleich, alle drei Werke erzielen jedoch
niedrigere Erlöse (1~P.). Im schlimmsten Fall entsteht ein \emph{ruinöser Preiskampf},
bei dem alle Anbieter Verluste erleiden. Sinnvoller ist ein Wettbewerb über Qualität und
Service (1~P.).

\textbf{d) Übernahme --- Angebotsmonopol} \hfill (3~P.: Erläuterung 1~P., je Grenze 1~P.)

Nach der Übernahme ist das Werk alleiniger Anbieter in der Region
(\emph{Angebotsmonopol}). Es muss keine Reaktionen von Konkurrenten mehr berücksichtigen
und kann den Preis als \emph{Preisfixierer} weitgehend selbst festsetzen. Grenzen der
Preissetzungsmacht sind z.\,B.:
\begin{itemize}
  \item die \emph{Nachfrage}: Mit steigendem Preis sinkt die Absatzmenge; Abnehmer
    verzichten auf Bauvorhaben oder weichen auf \emph{Substitute} aus (z.\,B. Holz-,
    Stahl- oder Ziegelbau);
  \item der Bezug aus \emph{entfernteren Werken bzw. Importe}, der sich bei stark
    steigenden Preisen trotz Transportkosten wieder lohnt;
  \item der drohende \emph{Markteintritt} neuer Anbieter, die von den hohen Preisen
    angelockt werden;
  \item die staatliche \emph{Wettbewerbsaufsicht} (Fusionskontrolle,
    Missbrauchsaufsicht durch das Bundeskartellamt --- vgl. Lektion~2).
\end{itemize}
\end{loesung}
